\documentclass[11pt,a4paper]{article}
\usepackage[margin=25mm,headheight=14pt]{geometry}
\usepackage{fontspec}
\setmainfont{lmroman10-regular.otf}[BoldFont=lmroman10-bold.otf,ItalicFont=lmroman10-italic.otf,BoldItalicFont=lmroman10-bolditalic.otf]
\setsansfont{lmsans10-regular.otf}
\setmonofont{lmmono10-regular.otf}[Scale=MatchLowercase]
\newfontfamily\authorsymbol{Noto Sans Symbols}[Scale=MatchUppercase]
\usepackage{amsmath,amssymb,amsthm,mathtools}
\usepackage{booktabs,longtable,array}
\usepackage{microtype}
\usepackage[hidelinks,unicode]{hyperref}
\usepackage{fancyhdr}
\pagestyle{fancy}
\fancyhf{}
\fancyhead[L]{\small\itshape A twelve-state composition tensor}
\fancyhead[R]{\small\thepage}
\renewcommand{\headrulewidth}{0pt}
\setlength{\parindent}{1.2em}
\setlength{\parskip}{0pt}
\setlength{\emergencystretch}{2em}
\newtheorem{theorem}{Theorem}
\newtheorem{proposition}[theorem]{Proposition}
\newtheorem{lemma}[theorem]{Lemma}
\newtheorem{corollary}[theorem]{Corollary}
\theoremstyle{definition}
\newtheorem{definition}[theorem]{Definition}
\DeclareMathOperator{\rank}{rank}
\DeclareMathOperator{\supp}{supp}
\newcommand{\R}{\mathbb R}
\newcommand{\Q}{\mathbb Q}
\newcommand{\ind}{\mathbf 1}
\hypersetup{pdftitle={Exact nonnegative rank of a twelve-state composition tensor},
 pdfauthor={lmtlssss🜐},pdfsubject={Mathematics: algebra and nonnegative tensor rank},
 pdfkeywords={nonnegative tensor rank, transformation monoid, rectangle partition, fooling set}}
\title{\LARGE Exact nonnegative rank of a\protect\\twelve-state composition tensor}
\author{lmtlssss{\authorsymbol 🜐}}
\date{26 October 2026}
\begin{document}
\maketitle
\thispagestyle{plain}
\begin{abstract}
We determine the nonnegative rank of the composition tensor of a specified
transformation monoid on four points. The monoid has twelve elements, and
its tensor has 144 nonzero coefficients. An explicit partition into 35
rectangles gives a nonnegative rank-one decomposition. A matching set of
35 incompatible entries proves the lower bound. The same entries also
specify determinant-one minors in the twelve output slices. Consequently,
the nonnegative tensor rank is exactly 35 over both the real and rational
numbers. The complete multiplication table, decomposition, lower-bound
certificates and an exact verifier are supplied.
\end{abstract}

\noindent\textit{Mathematics Subject Classification (2020).}
15A69; 15B48, 20M20.\par
\noindent\textit{Keywords.} Nonnegative tensor rank; transformation monoid;
rectangle partition; fooling set.

\section{The question}

A composition table records one output for each ordered pair of inputs.
The number of recorded products need not equal the number of rank-one
terms required to express the table. Here the distinction can be resolved
exactly: 144 products admit a decomposition into 35 nonnegative terms,
and every nonnegative decomposition requires at least 35 terms.

The object is the twelve-element monoid $S$ defined in Section~\ref{sec:monoid}.
Write its product as $m_i*m_j=m_{M(i,j)}$, and define
\begin{equation}\label{eq:tensor}
 C_{ijk}=\ind_{\{M(i,j)=k\}},\qquad 0\leq i,j,k<12.
\end{equation}
Thus $C\in\{0,1\}^{12\times12\times12}$. For each $k$, let $A_k$ be the
matrix with entries $(A_k)_{ij}=C_{ijk}$.

\begin{theorem}\label{thm:main}
The tensor $C$ has nonnegative rank $35$. Its output slices satisfy
\begin{equation}\label{eq:ranks}
\bigl(\rank_+(A_0),\ldots,\rank_+(A_{11})\bigr)
=(3,2,3,2,1,5,4,2,5,3,2,3).
\end{equation}
Each slice has the same ordinary matrix rank as its nonnegative rank.
An optimal decomposition of $C$ uses only $0$--$1$ factor vectors.
\end{theorem}

Nonnegative matrix rank is treated in Cohen and Rothblum~\cite{cohen1993};
nonnegative tensor rank is studied in Qi, Comon and Lim~\cite{qi2016}.
The rectangle-covering and fooling-set bounds used below are standard
support bounds; see Fiorini et al.~\cite{fiorini2013}. The contribution of
this note is the exact calculation and explicit certificates for $S$.

\section{The monoid}\label{sec:monoid}

Let $X=\{0,1,2,3\}$. A tuple $(a_0,a_1,a_2,a_3)$ denotes the map
$x\mapsto a_x$ on $X$. Consider the following twelve maps.

\begin{table}[ht]
\centering
\renewcommand{\arraystretch}{1.12}
\begin{tabular}{r l @ {\qquad\qquad} r l}
\toprule
$i$ & $m_i$ & $i$ & $m_i$\\
\midrule
0 & (0,0,0,0) & 6 & (1,1,2,2) \\
1 & (0,0,0,1) & 7 & (1,1,2,3) \\
2 & (0,1,1,1) & 8 & (2,2,2,2) \\
3 & (0,1,2,2) & 9 & (2,2,2,3) \\
4 & (0,1,2,3) & 10 & (2,3,3,3) \\
5 & (1,1,1,1) & 11 & (3,3,3,3) \\
\bottomrule
\end{tabular}
\caption{The elements of $S$, in lexicographic order.}\label{tab:maps}
\end{table}

We use the chronological convention
\begin{equation}\label{eq:composition}
 m_i*m_j=m_j\circ m_i.
\end{equation}
In particular, the row index acts first. The identity is $m_4$.
The six maps $m_1,m_2,m_3,m_7,m_9,m_{10}$ generate $S$ as a monoid.

For completeness, the other five nonidentity elements are obtained by
\[
\begin{aligned}
 m_0&=m_1*m_1,\quad m_5=m_1*m_7,\quad m_6=m_3*m_7,\\
 m_8&=m_9*m_3,\qquad m_{11}=m_9*m_{10}.
\end{aligned}
\]
The multiplication table in Appendix~\ref{app:table} shows closure.
Associativity follows from composition of functions, so these data define
the monoid without reference to an external presentation.

\section{Reduction to output slices}

For an entrywise nonnegative tensor $T$, its nonnegative rank is the least
$r$ for which
\begin{equation}\label{eq:rank}
 T=\sum_{\ell=1}^{r}u_\ell\otimes v_\ell\otimes w_\ell,
 \qquad u_\ell,v_\ell,w_\ell\geq0.
\end{equation}
The corresponding matrix definition uses sums of nonnegative rank-one
matrices. Zero summands may be removed.

The single-output condition in~\eqref{eq:tensor} has a useful consequence.

\begin{lemma}[Output-slice confinement]\label{lem:slices}
Let $f:I\times J\to K$ be a function between finite sets, and let
$T_{ijk}=\ind_{\{f(i,j)=k\}}$. Write $T_k=(T_{ijk})_{i,j}$. Then
\begin{equation}\label{eq:slice-sum}
 \rank_+(T)=\sum_{k\in K}\rank_+(T_k).
\end{equation}
\end{lemma}

\begin{proof}
Take a nonzero summand $u\otimes v\otimes w$ in a nonnegative decomposition
of $T$. Choose $i,j$ with $u_i v_j>0$. If $w$ were positive at two distinct
indices $k,k'$, the summand would be positive at both $(i,j,k)$ and
$(i,j,k')$. One of these entries is zero in $T$. Nonnegative summands
cannot cancel, so this is impossible. Hence $w$ is supported at a single
index $k$.

Every summand therefore belongs to one output slice. If $r_k$ summands
belong to slice $k$, they give a nonnegative matrix decomposition of $T_k$,
so $r_k\geq\rank_+(T_k)$. Summing gives the lower bound
in~\eqref{eq:slice-sum}. Conversely, choose a minimum nonnegative
decomposition of each $T_k$ and attach the coordinate vector $e_k$ to
each term. Their sum is a decomposition of $T$ with
$\sum_k\rank_+(T_k)$ terms.
\end{proof}

The lemma concerns the zero pattern of a function tensor. Associativity
is not needed. It also shows why a lower bound established separately
in each output slice can be added without assuming disjoint factors.

\section{The upper bound}

For nonempty $L,R\subseteq\{0,\ldots,11\}$, call $L\times R$ a
\emph{$k$-rectangle} when $M(i,j)=k$ for every $i\in L$, $j\in R$.
Let $\chi_L\in\{0,1\}^{12}$ denote the characteristic vector of $L$.
Such a rectangle gives the nonnegative rank-one tensor
$\chi_L\otimes\chi_R\otimes e_k$.

Table~\ref{tab:rectangles} lists 35 rectangles. Within each output slice
they are disjoint, and together they cover all 144 ordered input pairs.
Equivalently, the table gives the identity
\begin{equation}\label{eq:decomposition}
 C=\sum_{(k,L,R)\in\mathcal P}\chi_L\otimes\chi_R\otimes e_k,
 \qquad |\mathcal P|=35.
\end{equation}
Each coefficient can be read directly from the multiplication table:
for any $(i,j)$, exactly one listed rectangle contains that pair, and its
output label is $M(i,j)$.

For example, the rectangles with output $0$ are
\[
 \{0\}\times\{2,3,4\},\quad
 \{0,1,2,3,5,6,8\}\times\{0,1\},\quad
 \{4,7,9,10,11\}\times\{0\}.
\]
Their supports partition $A_0$. The other eleven slices are recorded in
the same form. The numbers of rectangles by slice are exactly the entries
of~\eqref{eq:ranks}. This proves the required upper bounds.

\section{The lower bound}

For a fixed output $k$, let $F_k$ be a set of positions $(i,j)$ with
$M(i,j)=k$. Call it a \emph{fooling set} if every two distinct positions
$(i,j),(u,v)\in F_k$ satisfy
\begin{equation}\label{eq:fooling}
 M(i,v)\ne k\quad\text{or}\quad M(u,j)\ne k.
\end{equation}

\begin{lemma}\label{lem:fooling}
If $F_k$ is a fooling set, then $\rank_+(A_k)\geq |F_k|$.
\end{lemma}

\begin{proof}
The positive support of a nonzero nonnegative rank-one matrix $uv^T$
is $\supp(u)\times\supp(v)$, a rectangle. A summand in a nonnegative
decomposition of $A_k$ must have its support inside the support of $A_k$.
If it covered both $(i,j)$ and $(u,v)$ from $F_k$, it would also be
positive at both crossed positions $(i,v)$ and $(u,j)$, contradicting
\eqref{eq:fooling}. Thus each summand covers at most one element of $F_k$.
Every element must be covered by some summand, proving the bound.
\end{proof}

The twelve fooling sets in Table~\ref{tab:witnesses} have sizes
\[
 3,2,3,2,1,5,4,2,5,3,2,3.
\]
For example, $F_0=\{(0,2),(1,1),(4,0)\}$. All three positions have output
$0$. The crossed entries $M(1,2)=1$, $M(4,2)=2$ and $M(4,1)=1$ establish
the three required incompatibilities. The remaining checks use the
same condition and the displayed multiplication table.

\begin{proof}[Proof of Theorem~\ref{thm:main}]
The rectangle partitions give $\rank_+(A_k)\leq r_k$, where $r_k$ is the
$k$th entry of~\eqref{eq:ranks}. Table~\ref{tab:witnesses} and
Lemma~\ref{lem:fooling} give $\rank_+(A_k)\geq r_k$.
Lemma~\ref{lem:slices} now yields
\[
 \rank_+(C)=\sum_{k=0}^{11}r_k=35.
\]
The decomposition~\eqref{eq:decomposition} uses $0$--$1$ vectors.

To obtain the ordinary slice ranks, retain the order of the pairs in
Table~\ref{tab:witnesses}. If $F_k=((i_1,j_1),\ldots,(i_{r_k},j_{r_k}))$,
form
\[
 B_k=\bigl(\ind_{\{M(i_s,j_t)=k\}}\bigr)_{1\leq s,t\leq r_k}.
\]
The row indices are distinct, as are the column indices. Direct evaluation
from Table~\ref{tab:multiplication} gives $\det B_k=1$ for every $k$.
Hence $r_k\leq\rank(A_k)\leq\rank_+(A_k)=r_k$.
\end{proof}

The determinant calculation gives a second lower-bound certificate. For
instance, the witnesses for output $5$ specify
\[
 B_5=\begin{pmatrix}
 1&1&0&0&0\\
 0&1&0&0&0\\
 1&1&1&1&0\\
 0&1&0&1&0\\
 0&1&0&1&1
 \end{pmatrix},\qquad \det B_5=1.
\]
Either the fooling-set checks or the minor determinants, together with
the rectangle partitions and Lemma~\ref{lem:slices}, proves the value 35.

\begin{corollary}\label{cor:rational}
The nonnegative rank of $C$ over $\Q$ is also 35.
\end{corollary}
\begin{proof}
The real lower bound applies to every rational decomposition, and
\eqref{eq:decomposition} has integer factor vectors.
\end{proof}

\section{Bilinear form and verification}

For $a,b\in\R^{12}$, define
\[
 p_k(a,b)=\sum_{M(i,j)=k}a_i b_j.
\]
The decomposition gives the polynomial identity
\begin{equation}\label{eq:bilinear}
 p_k(a,b)=\sum_{(k,L,R)\in\mathcal P}
       \left(\sum_{i\in L}a_i\right)
       \left(\sum_{j\in R}b_j\right).
\end{equation}
It uses 35 products of linear forms across all outputs. Among bilinear
decompositions with nonnegative coefficients in the two input factors
and the output factor, 35 is minimal. The identity itself holds for
signed inputs as well. The lower bound concerns nonnegative
factorizations; the unrestricted real and complex tensor ranks are
separate questions.

The supplementary file \texttt{certificate.json} contains the maps,
generators, multiplication table, rectangles and ordered witness pairs.
The program \texttt{verify.py} reconstructs the multiplication table from
the maps, checks the generated monoid, and verifies all 1,728 tensor
coefficients. It then checks every fooling-set condition and computes
the twelve witness determinants over the rational numbers. It uses
only the Python standard library. From this paper's directory, run
\begin{verbatim}
python3 verify.py
\end{verbatim}
The proof is the finite construction and lower bound displayed here.
The verifier checks their transcription and arithmetic; it requires
neither a search procedure nor an optimisation solver.

\appendix
\clearpage
\section{Multiplication table}\label{app:table}

Each entry is the index $M(i,j)$ determined by $m_j\circ m_i$.
The row acts first. In particular, row $4$ and column $4$ give the
identity operation.

\begin{table}[ht]
\centering
\small
\setlength{\tabcolsep}{5.1pt}
\renewcommand{\arraystretch}{1.22}
\begin{tabular}{r|rrrrrrrrrrrr}
\toprule
$i\backslash j$ & 0&1&2&3&4&5&6&7&8&9&10&11\\
\midrule
0 & 0 & 0 & 0 & 0 & 0 & 5 & 5 & 5 & 8 & 8 & 8 & 11 \\
1 & 0 & 0 & 1 & 1 & 1 & 5 & 5 & 5 & 8 & 8 & 9 & 11 \\
2 & 0 & 0 & 2 & 2 & 2 & 5 & 5 & 5 & 8 & 8 & 10 & 11 \\
3 & 0 & 0 & 2 & 3 & 3 & 5 & 6 & 6 & 8 & 8 & 10 & 11 \\
4 & 0 & 1 & 2 & 3 & 4 & 5 & 6 & 7 & 8 & 9 & 10 & 11 \\
5 & 0 & 0 & 5 & 5 & 5 & 5 & 5 & 5 & 8 & 8 & 11 & 11 \\
6 & 0 & 0 & 5 & 6 & 6 & 5 & 6 & 6 & 8 & 8 & 11 & 11 \\
7 & 0 & 1 & 5 & 6 & 7 & 5 & 6 & 7 & 8 & 9 & 11 & 11 \\
8 & 0 & 0 & 5 & 8 & 8 & 5 & 8 & 8 & 8 & 8 & 11 & 11 \\
9 & 0 & 1 & 5 & 8 & 9 & 5 & 8 & 9 & 8 & 9 & 11 & 11 \\
10 & 0 & 2 & 5 & 8 & 10 & 5 & 8 & 10 & 8 & 10 & 11 & 11 \\
11 & 0 & 5 & 5 & 8 & 11 & 5 & 8 & 11 & 8 & 11 & 11 & 11 \\
\bottomrule
\end{tabular}
\caption{The multiplication table of $S$.}\label{tab:multiplication}
\end{table}

\section{Lower-bound certificate}\label{app:lower}

Within each row of the following table, the pairs are ordered. This
order defines both the fooling set $F_k$ and the minor $B_k$ used in the
proof. The final two columns give its size and determinant.

\begin{table}[ht]
\centering
\small
\setlength{\tabcolsep}{6pt}
\renewcommand{\arraystretch}{1.2}
\begin{tabular}{r l r r}
\toprule
$k$ & Ordered pairs in $F_k$ & $r_k$ & $\det B_k$\\
\midrule
0 & $(0,2),\,(1,1),\,(4,0)$ & 3 & $1$ \\
1 & $(1,2),\,(4,1)$ & 2 & $1$ \\
2 & $(2,3),\,(3,2),\,(10,1)$ & 3 & $1$ \\
3 & $(3,4),\,(4,3)$ & 2 & $1$ \\
4 & $(4,4)$ & 1 & $1$ \\
5 & $(0,6),\,(3,5),\,(5,3),\,(6,2),\,(11,1)$ & 5 & $1$ \\
6 & $(3,7),\,(4,6),\,(6,4),\,(7,3)$ & 4 & $1$ \\
7 & $(4,7),\,(7,4)$ & 2 & $1$ \\
8 & $(0,10),\,(1,9),\,(4,8),\,(8,4),\,(9,3)$ & 5 & $1$ \\
9 & $(1,10),\,(4,9),\,(9,4)$ & 3 & $1$ \\
10 & $(2,10),\,(10,4)$ & 2 & $1$ \\
11 & $(0,11),\,(5,10),\,(11,4)$ & 3 & $1$ \\
\midrule
 & Total & 35 & \\
\bottomrule
\end{tabular}
\caption{Matching lower bounds for the twelve slices.}\label{tab:witnesses}
\end{table}

\clearpage
\section{The 35-term decomposition}\label{app:upper}

A row $(k,L,R)$ represents $\chi_L\otimes\chi_R\otimes e_k$.
The rows partition all 144 nonzero coefficients of $C$.

\small
\renewcommand{\arraystretch}{0.95}
\begin{longtable}{r l l}
\caption{The rectangle partition.}\label{tab:rectangles}\\
\toprule
$k$ & $L$ & $R$\\
\midrule
\endfirsthead
\multicolumn{3}{l}{\textit{Table \thetable\ continued.}}\\
\toprule
$k$ & $L$ & $R$\\
\midrule
\endhead
\bottomrule
\endfoot
0 & $\{0\}$ & $\{2,3,4\}$ \\
0 & $\{0,1,2,3,5,6,8\}$ & $\{0,1\}$ \\
0 & $\{4,7,9,10,11\}$ & $\{0\}$ \\
\addlinespace[0pt]
1 & $\{1\}$ & $\{2,3,4\}$ \\
1 & $\{4,7,9\}$ & $\{1\}$ \\
\addlinespace[0pt]
2 & $\{2\}$ & $\{2,3,4\}$ \\
2 & $\{3,4\}$ & $\{2\}$ \\
2 & $\{10\}$ & $\{1\}$ \\
\addlinespace[0pt]
3 & $\{3\}$ & $\{3,4\}$ \\
3 & $\{4\}$ & $\{3\}$ \\
\addlinespace[0pt]
4 & $\{4\}$ & $\{4\}$ \\
\addlinespace[0pt]
5 & $\{3,4\}$ & $\{5\}$ \\
5 & $\{5\}$ & $\{2,3,4\}$ \\
5 & $\{0,1,2,5\}$ & $\{5,6,7\}$ \\
5 & $\{11\}$ & $\{1\}$ \\
5 & $\{6,7,8,9,10,11\}$ & $\{2,5\}$ \\
\addlinespace[0pt]
6 & $\{4\}$ & $\{6\}$ \\
6 & $\{6\}$ & $\{3,4\}$ \\
6 & $\{3,6\}$ & $\{6,7\}$ \\
6 & $\{7\}$ & $\{3,6\}$ \\
\addlinespace[0pt]
7 & $\{7\}$ & $\{4\}$ \\
7 & $\{4,7\}$ & $\{7\}$ \\
\addlinespace[0pt]
8 & $\{0\}$ & $\{10\}$ \\
8 & $\{0,1,2,3,5,6\}$ & $\{8,9\}$ \\
8 & $\{4,7\}$ & $\{8\}$ \\
8 & $\{8\}$ & $\{4,7,9\}$ \\
8 & $\{8,9,10,11\}$ & $\{3,6,8\}$ \\
\addlinespace[0pt]
9 & $\{1\}$ & $\{10\}$ \\
9 & $\{9\}$ & $\{4,7\}$ \\
9 & $\{4,7,9\}$ & $\{9\}$ \\
\addlinespace[0pt]
10 & $\{2,3,4\}$ & $\{10\}$ \\
10 & $\{10\}$ & $\{4,7,9\}$ \\
\addlinespace[0pt]
11 & $\{0,1,2,3,4\}$ & $\{11\}$ \\
11 & $\{11\}$ & $\{4,7,9\}$ \\
11 & $\{5,6,7,8,9,10,11\}$ & $\{10,11\}$ \\
\end{longtable}
\normalsize

{\small
\begin{thebibliography}{9}
\bibitem{cohen1993}
J.~E.~Cohen and U.~G.~Rothblum,
Nonnegative ranks, decompositions, and factorizations of nonnegative matrices,
\textit{Linear Algebra and its Applications} \textbf{190} (1993), 149--168.
\href{https://doi.org/10.1016/0024-3795(93)90224-C}{doi:10.1016/0024-3795(93)90224-C}.

\bibitem{fiorini2013}
S.~Fiorini, V.~Kaibel, K.~Pashkovich and D.~O.~Theis,
Combinatorial bounds on nonnegative rank and extended formulations,
\textit{Discrete Mathematics} \textbf{313} (2013), no.~1, 67--83.
\href{https://doi.org/10.1016/j.disc.2012.09.015}{doi:10.1016/j.disc.2012.09.015}.

\bibitem{qi2016}
Y.~Qi, P.~Comon and L.-H.~Lim,
Semialgebraic geometry of nonnegative tensor rank,
\textit{SIAM Journal on Matrix Analysis and Applications} \textbf{37} (2016),
no.~4, 1556--1580.
\href{https://doi.org/10.1137/16M1063708}{doi:10.1137/16M1063708}.
\end{thebibliography}}
\end{document}
